% ==================================================
% CHAPTER 4: IMPLEMENTATION
% ==================================================

\chapter{Implementation}
\label{ch:implementation}

\section{Code Base}

The repository holds the backend (23 Python modules, about 4{,}000 lines), the web app (21 Svelte components), the Chrome extension, 551 tests, the evaluation scripts and the documentation. Table~\ref{tab:modules} lists the main backend modules.

\begin{table}[!htbp]
\centering
\small
\begin{tabularx}{\textwidth}{>{\ttfamily}p{0.25\textwidth} L}
\toprule
\textrm{\textbf{Module}} & \textbf{Responsibility} \\
\midrule
main.py & API endpoints, rate limits, log filtering, automatic-save re-check \\
auth.py, accounts.py & Token verification; linking sign-ins by verified email \\
library\_moves.py & Background merge of a library into the linked one \\
ingest.py & The save pipeline, under a per-user write lock \\
cleaning.py, chunking.py & Main-content extraction; overlapping chunks \\
embeddings.py & Model loading, batched passage and cached query embedding \\
retrieval.py & BM25 + KNN, fusion, page grouping, relevance floor \\
query\_intent.py & Site words, time phrases and URLs in queries \\
snippets.py & Sentence snippets for each result \\
solr\_client.py & All Solr access, with the user filter built in \\
page\_filter.py & Rules for pages automatic capture must refuse \\
summarizer.py & On-demand Gemini summaries with a time limit \\
reports.py & Stored problem reports and the email sender \\
youtube\_utils.py & Validated video ids and transcript processing \\
\bottomrule
\end{tabularx}
\caption{Main backend modules}
\label{tab:modules}
\end{table}

\section{Browser Extension}

The extension is a Manifest V3 extension \cite{chromemv3} with a popup, a
background service worker and scripts injected into pages only when needed.
It signs in through Cognito with \texttt{chrome.identity} and keeps its
tokens in \texttt{chrome.storage.local}, which web pages cannot read.

\subsection{Saving a Page}

\textbf{Remember this page} injects a content script that returns the page's
HTML, its visible text and its canonical URL. Single-page applications such
as Reddit and YouTube change the address without reloading and keep a stale
canonical tag, so a canonical URL is sent only when it describes the current
page (same host and path, and matching parameters). The popup confirms the
save and offers \textbf{Undo}.

\subsection{YouTube}

YouTube refuses transcript requests from servers. For a video, the
extension therefore fetches the caption track, title, channel, length and
description \emph{inside the user's own tab}, and sends them with the save
(ADR~0006). A video without captions is saved from its description.

\subsection{Automatic Capture}

Automatic capture is off by default and explained before it is turned on.
Rather than trying to list every site that should not be saved, it keeps a
page only if it passes three tests (swimlane diagram, Figure~\ref{fig:swimlane}):

\begin{enumerate}
  \item \textbf{The address is not private}: not email, banking, social
  media, a private application (documents, chat, cloud consoles, college
  portals), a local address, an account area, or a link carrying a secret.
  \item \textbf{The user read it}: the page was visible for 30 seconds, or
  scrolled at least halfway after 5 seconds. Pages bounced off are never
  saved.
  \item \textbf{It reads like content}: at least 150 readable words, not
  mostly links, no visible password field, no large editor or form, and not
  marked \texttt{noindex}.
\end{enumerate}

The same address rules are applied again by the server to every automatic
save, so even an outdated or faulty extension cannot store an inbox or a
login page. An automatic save never rewrites a page that is already in the
library; only a manual save updates one. The popup lists the day's automatic
saves with \textbf{Remove} and \textbf{Never this site}, and drops entries the
user has since deleted on the website.

\begin{figure}[!htbp]
\centering
\begin{subfigure}[b]{0.52\textwidth}
  \centering
  \screenshot[\textwidth]{figures/extension-popup.png}{extension popup}
  \caption{Popup after a save, with ``Auto-saved today''}
\end{subfigure}\hfill
\begin{subfigure}[b]{0.44\textwidth}
  \centering
  \screenshot[\textwidth]{figures/excluded-sites.png}{excluded sites page}
  \caption{Excluded sites (options page)}
\end{subfigure}
\caption{The Chrome extension}
\label{fig:extension}
\end{figure}

\section{Ingestion Pipeline}

All saves, from the API, the re-index scripts and the evaluation, go through
one function, \texttt{ingest\_page} (Figure~\ref{fig:seq-save}). Writes for one
user run under that user's lock, so two saves cannot race on the page
limit or leave duplicate chunks.

\subsection{URL Normalisation and Identity}

The same page can have many addresses. Normalisation lower-cases the scheme
and host, removes \texttt{www.}, fragments, trailing slashes and tracking
parameters (\texttt{utm\_*}, \texttt{fbclid}, \texttt{gclid}, \texttt{ref}),
and sorts the remaining parameters. The page id is

\[
\texttt{page\_id} = \mathrm{SHA256}(\texttt{user\_id} \,\|\, \texttt{"|"} \,\|\, \textit{normalised URL})_{[0:32]},
\]

so the same page saved twice updates one entry, and two users' copies of a
page can never collide.

\subsection{Cleaning and Chunking}

trafilatura \cite{barbaresi2021trafilatura} extracts the main text of the
HTML and drops navigation, footers and banners; if it finds too little, the
extension's visible text is used. Citation markers and junk lines are then
removed. The text is split into paragraphs and sentences and packed into
chunks of at most 150 words, each starting with about 30 words from the end
of the previous chunk, so an idea is never cut in half at a boundary. A
short final chunk merges into the one before.

\subsection{Embeddings}

Each chunk is embedded as ``title + chunk'' with \texttt{BAAI/bge-small-en-v1.5}
\cite{xiao2023bge} through Sentence Transformers \cite{reimers2019sbert}:
384-dimensional, L2-normalised vectors, computed in batches of 32 on the CPU.
Queries are embedded with BGE's retrieval instruction in front, and the last
1{,}024 query vectors are cached. A content hash and the model name are
stored with each page, so an unchanged page costs one lookup to re-save, and
a model change re-embeds pages the next time they are saved.

\section{Hybrid Retrieval}

Figures~\ref{fig:seq-search} and~\ref{fig:collab-search} show one search from query to results, and Figure~\ref{fig:dfd2} its data flow.

\subsection{Lexical and Semantic Retrieval}

BM25 \cite{robertson2009bm25} runs through Solr's edismax parser over
\texttt{title\^{}2 text}, with a phrase boost so pages containing the query as
a phrase rank higher, and with field syntax in user input disabled. The
query vector is searched with Solr's HNSW KNN \cite{solr}, using the user
filter as a \emph{pre-filter} so all 50 neighbours belong to the caller. The
two run concurrently.

\subsection{Site, Time and URL in the Query}

People often remember \emph{where} or \emph{when} they saw something. The
query-intent module recognises:

\begin{itemize}
  \item \textbf{site words}: names derived from the domains the user has
  actually saved (\texttt{en.wikipedia.org} $\rightarrow$ \textit{wikipedia},
  \texttt{docs.nvidia.com} $\rightarrow$ \textit{nvidia}), joined words
  (\textit{rotten tomatoes}) and a few nicknames (\textit{yt}, \textit{gh});
  \item \textbf{time phrases}: \textit{yesterday}, \textit{3 weeks ago},
  \textit{last month}, \textit{in march}, \textit{31 august}, \textit{recently};
  \item \textbf{a URL or domain} as the whole query.
\end{itemize}

When one is found, the rest of the query is searched again within those
domains and dates, and the extra lists join the fusion as boosts rather than
filters, so a better page from elsewhere still wins.

\subsection{Reciprocal Rank Fusion}

BM25 scores and cosine similarities live on different scales. The prototype
min-max normalised and averaged them, which lets one outlier shift every
other score. The final system fuses \emph{ranks} with Reciprocal Rank Fusion
\cite{cormack2009rrf} (ADR~0002):

\[
\mathrm{RRF}(d) \;=\; \sum_{\ell \,\in\, \text{lists}} \frac{w_\ell}{k + r_\ell(d)},
\qquad k = 60,\; w_\ell = 1,
\]

where $r_\ell(d)$ is the rank of chunk $d$ in list $\ell$, and a chunk
missing from a list gets nothing from it. Chunks found by both retrievers
rise to the top.

\subsection{Page Grouping and the Relevance Floor}

Users see pages, not chunks. Each page is scored by its best chunk,
$S_{\text{page}} = s_{\text{best}}$; Section~\ref{sec:bonus} explains why an
earlier bonus for a second matching chunk was removed. KNN always returns
neighbours, however unrelated, so a page found only by meaning is kept only
if its best cosine similarity is at least $0.55$ (ADR~0003). A page found by
keywords must hold at least 75\% of the query's words, or else meet the same
similarity bar: BM25 accepts a page that shares any single word, and for
``greatest football player of all time'' that would otherwise admit the
Nintendo article (``the greatest games of all time'').

\subsection{Snippets}

For each result, two or three consecutive readable sentences around the
sentence with the most query terms are shown, with the terms highlighted.
When no sentence shares a term (a search by meaning), the first readable
sentences are compared with the query by embedding instead.

\section{Summaries}

\textbf{Summarise} on a result asks Google Gemini \cite{gemini} for a short
summary of that page, adapted to the query, with a 20-second limit. Summaries
are on demand only, so search never waits on an external model, and the
system works fully without a Gemini key (ADR~0008).

\section{Web Application}

The web app is a Svelte \cite{svelte} single-page application served by
Caddy. Figure~\ref{fig:webapp} shows its main screens.

\begin{figure}[!htbp]
\centering
\begin{subfigure}[b]{\textwidth}
  \centering
  \screenshot[0.92\textwidth]{figures/search-results.png}{search results}
  \caption{A search naming a site (``nvidia docs on cuda graphs''): the matching
  snippet with query words highlighted, and the page's summary opened}
\end{subfigure}

\vspace{4mm}
\begin{subfigure}[b]{\textwidth}
  \centering
  \screenshot[0.7\textwidth]{figures/library.png}{library}
  \caption{The library, newest first (a delete button appears on each row on hover)}
\end{subfigure}
\caption{The web application}
\label{fig:webapp}
\end{figure}

Deleting is optimistic: the page disappears at once and a toast offers
\textbf{Undo} for six seconds before the request is sent. If the tab is
closed or reloaded in that window, the delete is sent immediately and the
reloaded page keeps it hidden. Voice search uses the browser's speech
recognition to fill the search box.

\section{Problem Reports}

\textbf{Report a problem}, on the website and in the extension, shows the
user exactly what will be sent (message, app version, page, browser, last
error) and stores the report in SQLite before a background sender emails it
through Amazon SES with retries (ADR~0007).

\section{Documentation Site}

The documentation is a 48-page MkDocs Material \cite{mkdocsmaterial} site
published on GitHub Pages, with a user guide, architecture and data-flow
diagrams, the API reference, the evaluation, the decision records and a
glossary. Its search box finds sections \emph{by meaning} as well as by
keywords, using the same idea as MemoryLane itself at a smaller scale: at
build time every section is embedded with the static model Model2Vec
\texttt{potion-base-8M} \cite{model2vec}, and in the browser a query is
tokenised, averaged into a vector and compared with all sections in a few
milliseconds, with no neural network running. The results are fused with the
keyword results by RRF. The files (about 8\,MB) load only when the search box
is first focused.
